% -----------------------------------------------------------------------------
% This file is automatically generated.
% Do not edit manually.
% Source: notebooks/support/hypothesis-testing/support.Rmd
% -----------------------------------------------------------------------------


\noindent We select the sample with the selection method \ttblue{selMeth = "random"} and order it in decreasing order, displaying the six largest book values.

\par\addvspace{\topsep}
\begingroup
\fvset{listparameters={%
  \setlength{\topsep}{0pt}%
  \setlength{\partopsep}{0pt}%
  \setlength{\parsep}{0pt}%
  \setlength{\itemsep}{0pt}%
}}
\begin{Verbatim}[breaklines=true,formatcom=\color{ada_blue}]
multiple.select(selMeth="random", seed=1)
sample = multiple.sample
sample.sort_values("bv", ascending=False).head()
\end{Verbatim}
\begin{Verbatim}[breaklines=true,formatcom=\color{ada_light_blue}]
          item       bv      cum     unit previous
938  201719763 12049.70  6248492 10416.59  6236442
1070 201719763 12049.70  6248492  5378.04  6236442
243  201734940 11981.19 12301198   782.72 12289216
543  201723585 10377.48  8822377  8218.39  8811999
501  201710344  9918.25  7486492  4816.81  7476574
822  201710344  9918.25  7486492  6251.71  7476574
\end{Verbatim}
\endgroup
\par\addvspace{\topsep}
\noindent This example demonstrates that invoices 201719763 and 201710344 were selected twice.
